\documentclass[10pt,conference]{IEEEtran}

\usepackage{amsmath,amssymb}
\usepackage{graphicx}
\usepackage{booktabs}
\usepackage{cite}
\usepackage{url}

\begin{document}

\title{A Beam-Conditioned 3D U-Net for Multi-Modality Radiotherapy Dose\\
Prediction, with a Physics-Based Extension to Proton Therapy:\\
System Design and Results for a Public Radiotherapy Dose-Prediction\\
Grand Challenge}

\author{
    \IEEEauthorblockN{Author name withheld for double-blind review}
    \IEEEauthorblockA{Affiliation withheld for double-blind review}
}

\maketitle

% =====================================================================
\begin{abstract}
Monte Carlo (MC) dose calculation is the accuracy gold standard in
radiotherapy treatment planning but is too slow for interactive or
real-time use, motivating learned surrogate models that approximate MC
dose distributions directly from patient imaging and beam geometry. We
describe our entry to a public radiotherapy dose-prediction grand
challenge across two
completed tasks -- Task 1 (photon dose prediction on CT) and Task 2
(photon dose prediction on MRI) against Geant4 MC ground truth -- and a
third, in-progress extension to proton pencil-beam-scanning therapy
(Task 3). Our system combines a ray-traced photon beam encoder, which
converts gantry angle, isocenter, and multi-leaf collimator (MLC)
aperture into a 3D beam-path feature map, with a 3D U-Net conditioned on
that map to predict per-control-point dose directly on the patient
imaging grid. We train with a composite loss aligned with the
challenge's official Masked MAE metric, augmented with two auxiliary
full-field terms; we present a case study in metric-driven debugging in
which a systematic leakage of predicted dose mass outside the patient
body -- invisible to our primary training metric but catastrophic to two
of the challenge's other four scored metrics -- was diagnosed, verified
empirically against ground-truth data, and corrected. Task 1 and Task 2
models are trained, evaluated on all four locally computable official
metrics, and packaged into Grand-Challenge-compliant inference
containers, each independently validated end-to-end against a real
patient input. For Task 3, we extend the framework to modality-agnostic
proton pencil-beam physics -- an analytic single-Gaussian beamlet model
derived from the Bragg-Kleeman range-energy relation, Bortfeld range
straggling, and the Highland multiple-Coulomb-scattering formula --
report a GPU memory architecture fix that was necessary to make training
at proton's substantially larger voxel grid feasible at all, and present
training-loss trends from an in-progress run without claiming final
validated results. We report all four tasks' status honestly, including
open risks (an MR-container inference-latency measurement that has not
yet been decomposed per beam, and Task 2's early-epoch overfitting), and
the concrete work remaining before final submission.
\end{abstract}

\begin{IEEEkeywords}
dose prediction, deep learning, radiotherapy, 3D U-Net, Monte Carlo
surrogate, grand challenge, photon dose calculation, proton therapy,
pencil beam scanning
\end{IEEEkeywords}

% =====================================================================
\section{Introduction}

Accurate radiotherapy treatment planning requires computing the 3D
radiation dose a patient will absorb for a given beam arrangement.
Monte Carlo (MC) simulation, which stochastically tracks individual
particle histories through a patient's imaging-derived geometry, is the
accepted accuracy standard for this calculation but is computationally
expensive -- a single plan can take minutes to hours to converge to low
statistical noise. This cost restricts how many plan variants a
clinician can practically evaluate, motivating machine-learned surrogate
models that approximate MC dose output directly, at a fraction of the
runtime, so that dose feedback can approach real time.

This grand challenge poses the task as a supervised learning
problem across four task variants spanning two radiation modalities
(photon and proton) and two imaging modalities (CT and MRI): given a
patient volume and a treatment beam's geometric parameters, predict the
3D dose distribution that a reference MC engine (Geant4) would compute,
at the granularity of individual \emph{control points} for photon
therapy or individual \emph{beamlets} for proton pencil-beam scanning.
Submissions are scored on a five-metric suite spanning per-beam accuracy
in the high-dose region, depth-dose-curve shape fidelity, full-plan dose
accuracy stratified by dose level, a clinically-motivated 3D Gamma
Index, and a DVH-based clinical score -- and are further required to run
within a strict per-beam inference budget on the challenge's reference
GPU hardware.

This paper documents the system we built across three of the four task
variants. Tasks 1 (photon/CT) and 2 (photon/MR) are complete: both are
trained, evaluated against all four locally computable official metrics
on held-out patients, and packaged into tested, Grand-Challenge-compliant
inference containers. We use Task 1's first full training run as a case
study in metric-driven debugging: it produced deceptively good numbers
on one metric while failing badly on others, and we walk through how we
diagnosed the root cause, verified our hypothesis against real data
before touching code, and corrected it -- a fix that then transferred
directly to Task 2 by construction, since both tasks share the same
model and loss code. Task 3 (proton/CT) extends the same architecture to
a physically different beam model; we describe the analytic pencil-beam
physics we derived and implemented for it, a GPU memory architecture
fix required to make training tractable on proton's larger voxel grid,
and report the training run's current status without presenting
unvalidated numbers as final results. Task 4 (proton/MR) has not yet been
started and is discussed only as remaining work.

% =====================================================================
\section{Problem Formulation}

\subsection{Photon Tasks (1: CT, 2: MR)}

Tasks 1 and 2 provide, per patient, a planning volume (CT or MRI), a set
of photon beams, and for each beam a sequence of control points
describing the beam's geometry at that delivery instant (gantry angle,
isocenter, and MLC leaf positions for an Elekta Versa HD linear
accelerator with an Agility 160-leaf MLC). Ground-truth dose is provided
per control point as a 3D volume on the same grid as the source imaging,
computed by Geant4 MC simulation. The learning task is to predict this
per-control-point 3D dose volume from the imaging volume and that
control point's beam parameters alone. Each (patient, beam, control
point) triple is treated as an independent training example. The two
tasks differ only in imaging modality and its associated preprocessing
(Section~\ref{sec:preproc}); the beam geometry, control-point structure,
and dose grid conventions are identical, which is why a single
architecture and training pipeline serves both.

\subsection{Proton Task (3: CT)}

Task 3 poses a physically different problem. Rather than control points
along a continuous gantry rotation, a proton plan is described as a set
of \emph{beams}, each containing multiple \emph{rays} (fixed source and
target positions), each ray containing multiple \emph{beamlets} at
discrete energies -- pencil-beam scanning (PBS) delivers dose as a
superposition of many narrow, mono-energetic proton beamlets rather than
a broad, MLC-shaped photon field. Ground-truth dose is again provided
per beamlet as a 3D volume, but on a substantially finer, patient-specific
grid ($1\times1\times3$\,mm spacing, with in-plane extents up to
$\sim$500 voxels, versus the fixed $2\times2\times2$\,mm,
$256^3$-voxel grid used for Tasks 1--2). The learning task is again to
predict a per-beamlet dose volume, this time from a proton beamlet's
source, target, and energy.

% =====================================================================
\section{Method}

\subsection{Data Preprocessing}
\label{sec:preproc}

Each volume is resampled to the modality's reference resolution (an
isotropic $2\times2\times2$\,mm grid for Tasks 1--2), then center
cropped or padded to a fixed shape so that all examples share a common
tensor shape for batching ($256^3$ for Tasks 1--2; $176\times512\times512$
at $3\times1\times1$\,mm spacing for Task 3, matching the finer proton
dose grid). CT intensities are normalized directly; MR intensities are
normalized with a separate, modality-appropriate transform, since MR
lacks CT's absolute Hounsfield-unit scale. A body mask is derived
directly from the source imaging in both cases (a CT-appropriate
threshold for Task 1, an MR-appropriate one for Task 2), used both as an
auxiliary training signal (Section~\ref{sec:loss}) and, for Task 3, as
the reference surface a proton beamlet's entry depth is measured from
(Section~\ref{sec:proton-encoder}). Ground-truth dose values are scaled
by a fixed constant, chosen so that typical per-control-point dose
maxima (empirically $\sim5\times10^{-5}$ to $1.3\times10^{-4}$\,Gy,
measured directly from 120 real dose files) land near $O(1)$, so that
regression targets sit in a numerically well-conditioned range for the
network's unconstrained output.

\subsection{Photon Beam Geometry Encoding}
\label{sec:photon-encoder}

The photon beam encoder converts a control point's geometric parameters
-- gantry angle, isocenter position, and per-leaf MLC aperture
boundaries -- into a dense 3D feature map aligned with the imaging grid,
via geometric ray tracing: for each voxel, the encoder determines
whether that voxel lies within the beam's rectangular field as
delimited by the current MLC leaf positions, projected along the beam's
central axis at the given gantry angle. This produces a soft-edged
beam-path mask volume that is concatenated with the source imaging as a
second input channel to the network. The coordinate grid used for ray
tracing depends only on volume shape and voxel spacing, both of which
are fixed identically across every patient after the preprocessing in
Section~\ref{sec:preproc}; we therefore build this grid once and share
it across all patients, applying each patient's isocenter as a cheap
per-sample offset rather than rebuilding the grid per sample. Combined
with GPU execution of the ray-tracing step, this gave an approximately
78$\times$ speedup over a naive per-sample CPU implementation, which was
necessary to keep both training throughput and per-beam inference
latency within budget.

\subsection{Proton Beam Geometry Encoding}
\label{sec:proton-encoder}

Task 3's beamlet geometry -- a source, a target, and a single energy --
admits no MLC-aperture analogue, so we designed an analytic
pencil-beam encoder from first-principles proton physics rather than
reusing the photon ray-tracing approach. For a beamlet of energy $E$
(MeV), we compute:
\begin{itemize}
    \item \textbf{Range} via the Bragg-Kleeman relation,
        $R\,\text{[mm]} = 10 \cdot 0.0022 \cdot E\,\text{[MeV]}^{1.77}$;
    \item \textbf{Range straggling} via Bortfeld's mono-energetic
        approximation~\cite{bortfeld1997},
        $\sigma_R\,\text{[mm]} = 10 \cdot 0.012 \cdot R\,\text{[cm]}^{0.935}$;
    \item \textbf{Lateral spread} via the Highland multiple-Coulomb-scattering
        formula~\cite{highland1975}, evaluated in water
        ($X_0 = 360.8$\,mm) using relativistic beam kinematics for a
        proton (rest mass 938.272\,MeV) at energy $E$.
\end{itemize}
For each voxel, we compute its signed depth along the beamlet axis and
its lateral distance from that axis, then evaluate a single Gaussian in
depth (centered at the entry-to-body depth plus range, width
$\sigma_R$) and a single Gaussian in the lateral direction (width given
by the Highland formula), matching the single-Gaussian pencil-beam
approximation used in the challenge's own reference beam model.
Entry-to-body depth is found per ray by marching the depth axis against
the body mask from Section~\ref{sec:preproc}, since a beamlet's source
point sits in air outside the patient rather than at the body surface --
an early version of this encoder measured range from the source instead
of the body entry point and produced a systematically wrong (all-zero,
after range-gating) mask; we caught this by validating peak-dose-location
distance against real beamlet dose files (3--17\,mm accuracy across five
sampled beamlets) before trusting the encoder in training.

\subsubsection{A shared-coordinate-grid memory fix, discovered by Task 3}

Both encoders were originally designed to cache one full-volume
coordinate-grid tensor per \emph{patient}, since coordinate grids are
expensive enough to build that avoiding a rebuild per sample was
essential for training throughput (Section~\ref{sec:photon-encoder}).
This was safe for Task 1--2's fixed $256^3$ grid ($\sim$200\,MB per
cached tensor, $\sim$13.4\,GB across 67 train patients) but fatal for
Task 3's much larger $176\times512\times512$ grid ($\sim$553\,MB per
tensor, $\sim$37\,GB across the same number of patients) -- Task 3
training crashed with a CUDA out-of-memory error partway through its
first epoch, after enough distinct shuffled patients had been visited to
fill the cache. We root-caused this via the memory arithmetic above
rather than treating it as a generic "reduce batch size" problem, and
fixed it by recognizing that the coordinate grid itself is
patient-\emph{independent} once volume shape and voxel spacing are
fixed by preprocessing (Section~\ref{sec:preproc}) -- only the
per-patient isocenter/origin offset actually varies, and that is a
cheap $(3,)$ tensor rather than a full volume. Refactoring both encoders
to build one shared, origin-free coordinate grid and apply the
per-patient origin as a lightweight offset collapsed the cache to a
single shared entry, eliminating the OOM without changing model output
(re-validated: identical beamlet peak-distance numbers before and after
the refactor).

\subsection{Model Architecture}

We use a 3D U-Net (five-level encoder/decoder, channel widths
$16$--$32$--$64$--$128$--$256$, stride-2 downsampling at each level, two
residual units per block) taking the two-channel (source imaging, beam
mask) volume as input and predicting a single-channel dose volume of the
same spatial size. Beam conditioning is early-fused: the beam mask is
treated as a second input channel alongside the imaging volume rather
than injected at a later layer, so that the network's earliest
convolutional features already have access to beam geometry. The
identical architecture is used, unmodified, for all three tasks; only
the input preprocessing and beam encoder differ.

\subsection{Training Objective}
\label{sec:loss}

Our loss function is designed around the challenge's own primary
per-beam metric, \emph{Masked MAE}: mean absolute error restricted to
voxels in the high-dose region (dose $\geq 10\%$ of that beam's maximum
dose), normalized by that maximum. We train a directly differentiable
version of this exact quantity rather than a loose proxy, so that the
training objective and the scored metric are aligned as closely as
possible.

Optimizing Masked MAE alone gives the network no gradient signal outside
the high-dose region, which risked the network treating everything
outside that mask as unconstrained -- fine for Masked MAE itself, but
damaging to metrics that examine the full dose distribution (the Gamma
Index and DVH-based clinical score in particular). We therefore add an
auxiliary full-field term restricted to voxels inside the patient body
mask, giving the network gradient signal in the low- and mid-dose region
as well.

\paragraph{An out-of-body dose leakage bug, diagnosed by metric, not by inspection}
After Task 1's first full training run (10 epochs), Masked MAE looked
reasonable ($\approx 9.5\%$), but two of the other four official
metrics -- IDD curve distance and Gamma Index pass rate -- were badly
degraded, in a way inconsistent with a model that was simply
undertrained (Table~\ref{tab:t1-fix}). We hypothesized that the two loss
terms above, which only ever examine voxels \emph{inside} the high-dose
mask or the body mask respectively, gave the network no reason to
suppress dose predicted \emph{outside} the patient's body -- a region
that never contributes to the training loss regardless of how much
(incorrect) dose mass the network places there, but does contribute to
metrics that look at the full 3D distribution or an along-axis
depth-dose curve. Before changing any code, we verified this hypothesis
directly against real validation data: ground-truth dose outside the
body mask measured $\approx 1\times10^{-8}$ on average across sampled
control points, five to six orders of magnitude below typical in-body
dose maxima ($\approx 10^{-4}$), confirming that dose outside the body
is expected to be physically negligible and that any systematic leakage
there is a genuine model defect rather than a property of the data.

We corrected this by adding a third loss term that penalizes the
model's mean absolute error against (near-zero) ground truth
specifically in the region outside the body mask, normalized the same
way as the other two terms. The combined objective is
\begin{equation}
    \mathcal{L} = w_m \, \mathcal{L}_{\text{masked}}
                + w_b \, \mathcal{L}_{\text{body}}
                + w_o \, \mathcal{L}_{\text{outside}},
\end{equation}
with $w_m = 1.0$ and $w_b = w_o = 0.2$ -- the metric-aligned term
dominates, with the other two acting as smaller regularizers. Because
this fix lives entirely in shared loss/training code, it applied to
Task 2 automatically without any Task-2-specific change.

\subsection{Training Procedure}

We optimize with AdamW (learning rate $2\times10^{-4}$, weight decay
$1\times10^{-5}$) and a cosine-annealing learning-rate schedule, using
automatic mixed precision and gradient-norm clipping (max norm 1.0) for
stability. Batch size is fixed at 1 by design (each example requires its
own beam-encoder call regardless of batching, since encoders are not
vectorized across samples with different origins), with gradient
accumulation available for a larger effective batch. Patient identity,
not control-point/beamlet identity, defines the train/validation split
(90/10), so that validation measures generalization to unseen patients.
A finite-value check before both the backward pass and the optimizer
step skips (rather than applies) any step producing non-finite loss or
gradients, since an earlier unguarded run silently degraded to a
constant NaN output partway through its first epoch, undetected for the
remainder of a multi-day run.

\subsection{Evaluation Protocol}

We evaluate against all four of the challenge's locally quantifiable
official metrics on held-out patients: Masked MAE (per beam); IDD
(Integrated Depth-Dose) curve distance, the root-mean-square deviation
between predicted and ground-truth dose profiles along the beam's
central axis, normalized by peak ground-truth IDD; Stratified Plan MAE,
a full-plan (all of a patient's beam doses summed) L1 error computed
separately within high/mid/low dose bands relative to a
prescription-dose reference; and 3D local Gamma Index at the clinically
strict 1\%/1\,mm criterion, using each plan's own maximum dose as a
prescription-dose stand-in. The fifth official metric, a DVH-based
clinical score, requires structure/contour data (PTV and organ-at-risk
delineations) not present in the locally available training data; we
report it as unavailable rather than approximate it with a substitute
that could misrepresent clinical relevance.

\section{Results}
\label{sec:results}

\subsection{Task 1: Photon/CT}

Table~\ref{tab:t1-fix} reports validation-set metrics (8 held-out
patients, 540 control points each) before and after the out-of-body
correction of Section~\ref{sec:loss}, both evaluated from a fine-tune
resumed at the pre-fix run's best checkpoint. The fix was confirmed
effective \emph{before} the full fine-tune completed, by evaluating the
first post-fix epoch alone.

\begin{table}[!t]
\centering
\caption{Task 1 (Photon/CT) validation metrics, pre- vs. post-fix}
\label{tab:t1-fix}
\begin{tabular}{@{}lcc@{}}
\toprule
Metric & Pre-fix & Post-fix \\
\midrule
Masked MAE            & 0.0953 (9.5\%) & 0.0949 (9.5\%) \\
IDD curve distance     & 126.5          & \textbf{0.56} \\
Stratified Plan MAE   & 0.0716 (7.2\%) & 0.0509 (5.1\%) \\
Gamma pass rate (1\%/1\,mm) & 0.0135 (1.3\%) & 0.0328 (3.3\%) \\
\bottomrule
\end{tabular}
\end{table}

IDD curve distance improved roughly $225\times$ and now sits solidly in
the expected $O(10^{-2}\text{--}10^{0})$ range; Stratified Plan MAE and
Gamma pass rate both improved substantially as well, while Masked MAE
(the only metric the pre-fix model was already directly optimizing
inside the high-dose mask) was essentially unchanged, exactly as
expected for a fix that targets an out-of-mask failure mode. Gamma pass
rate remains low in absolute terms; we believe this is at least
partially an artifact of our local prescription-dose stand-in (each
plan's own maximum dose, in the absence of real RTSTRUCT prescription
data locally), rather than purely a model-accuracy issue, since the
Gamma Index computation is sensitive to this reference value -- but we
have not yet been able to confirm this against the real challenge test
set's prescription data. The full 10-epoch fine-tune completed without
further crashes; validation loss was lowest at the first post-fix
epoch, with mild re-overfitting in later epochs, so the reported
checkpoint is that first post-fix epoch, not the final one.

\subsection{Task 2: Photon/MR}

Table~\ref{tab:t2} reports Task 2 validation metrics (same 8 held-out
patients, 540 control points each) using the best checkpoint by
validation loss.

\begin{table}[!t]
\centering
\caption{Task 2 (Photon/MR) validation metrics}
\label{tab:t2}
\begin{tabular}{@{}lc@{}}
\toprule
Metric & Value \\
\midrule
Masked MAE                    & 0.1025 (10.3\%) \\
IDD curve distance             & 4.32 \\
Stratified Plan MAE           & 0.0920 (9.2\%) \\
Gamma pass rate (1\%/1\,mm)   & 0.0309 (3.1\%) \\
\bottomrule
\end{tabular}
\end{table}

The best Task 2 checkpoint occurred at the first training epoch;
validation loss degraded steadily in every subsequent epoch while
training loss continued to improve smoothly, a clear overfitting
signature we detected by comparing the two curves rather than trusting
training loss alone, and which we handle correctly by construction
since our checkpointing logic already always saves the
best-by-validation-loss checkpoint separately from the most recent one.
Gamma pass rate (3.1\%) is close to Task 1's post-fix value (3.3\%),
consistent with the same local prescription-approximation limitation
discussed above rather than an MR-specific problem; IDD curve distance
is higher than Task 1's post-fix value but still in a qualitatively
reasonable range, and Masked MAE and Stratified Plan MAE are both modestly
higher than Task 1, plausibly reflecting the shorter effective training
(one well-generalizing epoch, versus Task 1's full fine-tune) rather
than an inherent MR-versus-CT accuracy gap -- distinguishing between
these two explanations is listed as remaining work
(Section~\ref{sec:remaining}).

\subsection{Inference Speed}

Task 1's measured per-control-point inference latency inside the
packaged submission container, after profiling and removing three
avoidable costs (a redundant CPU array copy, a duplicate SimpleITK
round-trip for cutoff thresholding, and unwarmed CUDA/cuDNN kernel
selection on the first real call), is $\approx$0.16\,s/control point on
our development GPU (NVIDIA A100) -- roughly $6\times$ under the
challenge's 1-second-per-beam limit. Task 2's container was validated
end-to-end against a real patient input (health check, warmup, and a
successful \texttt{/invoke} call producing correct output), completing
in 1.94\,s total for that call; unlike Task 1, this has not yet been
decomposed into a controlled per-control-point measurement, so whether
Task 2 meets the same per-beam budget with the same margin is an open
question we flag explicitly rather than assume by analogy to Task 1.
Timing on the challenge's actual reference hardware (an AWS g5 instance)
has not yet been independently confirmed for either task.

\subsection{Submission Packaging}

Both Task 1 and Task 2 are packaged behind the challenge's required HTTP
\texttt{invoke} API (a \texttt{/health} endpoint the platform polls
until the model is loaded and warmed, and a \texttt{/invoke} endpoint
that executes one prediction call) inside independent Docker containers,
following the challenge's documented input/output contract: imaging and
beam-metadata inputs are provided in fixed, always-fully-populated
10-slot batches (unused slots contain $1\times1\times1$ placeholder
volumes), and predicted dose is written back in the same 10-slot
structure, stacked into genuine 4D volumes via SimpleITK's
\texttt{JoinSeries} rather than a raw array stack (which the challenge
documentation notes can silently corrupt multi-beam output), on a grid
that exactly matches each input volume's size, spacing, origin, and
direction. Both containers were validated end-to-end locally against
real patient input, confirming correct startup, correct output geometry
and content, and (for Task 1) per-beam latency well under budget.

\section{Extension to Proton Therapy (Task 3)}

Task 3 reuses the architecture, loss, and training procedure of Tasks
1--2 unmodified, substituting the proton pencil-beam encoder of
Section~\ref{sec:proton-encoder} and the finer $176\times512\times512$
grid preprocessing of Section~\ref{sec:preproc}. The full 423\,GB
proton training dataset (spanning $\sim$77 patients and $\sim$81{,}000
beamlets, matching the challenge's documented scale) was downloaded and
verified against real patient geometry (grid spacing, extent, and beam
JSON schema cross-checked on five real patients) before training began.
Training is ongoing at the time of writing. We report training-loss
trends only, and explicitly do not present a validated held-out result
for this task: the first training epoch completed with a validation
loss of 0.1032, and a second epoch is in progress; given Task 2's
demonstrated overfitting pattern (best result at the very first epoch),
our stopping criterion is to compare each subsequent epoch's validation
loss against the running best and stop once it degrades, rather than
committing in advance to a fixed epoch count. Task 3 evaluation code
(extending \texttt{evaluate.py} to the proton beamlet index/encoder,
analogous to but not identical in structure to the photon evaluation
path) and submission packaging have not yet been built.

\section{Remaining Work}
\label{sec:remaining}

The following items remain, in priority order:

\begin{enumerate}
    \item \textbf{Confirm Task 2's per-control-point inference latency}
        against the 1-second-per-beam budget with the same rigor applied
        to Task 1, rather than relying on a single aggregate
        \texttt{/invoke} timing.
    \item \textbf{Upload the completed Task 1 and Task 2 containers} to
        the Grand Challenge platform -- the platform's actual
        build/run/scoring behavior has not yet been exercised end-to-end
        with our real account, and is currently the largest untested
        risk for both completed tasks.
    \item \textbf{Complete or early-stop Task 3 training}, run its full
        evaluation suite, and determine whether Task 3's
        (currently unconfirmed) accuracy and inference latency meet the
        same bar as Tasks 1--2 before committing engineering time to its
        submission container.
    \item \textbf{Build Task 3's submission container}, which will
        require different input parsing (beam/ray/beamlet JSON rather
        than control-point JSON) than the shared Task 1--2 container
        pattern.
    \item \textbf{Task 4 (proton/MR)} has not yet been started; it would
        reuse Task 3's proton physics with Task 2's MR preprocessing, by
        the same composition principle already demonstrated between
        Tasks 1 and 2.
    \item \textbf{Investigate the Gamma Index pass-rate gap} between our
        local prescription-stand-in evaluation and the real challenge
        test set's prescription-based scoring, to determine how much of
        the currently low pass rate (Tasks 1--2, 3.1--3.3\%) reflects
        genuine model inaccuracy versus an artifact of the local
        evaluation approximation.
\end{enumerate}

% =====================================================================
\section{Conclusion}

We have described a beam-conditioned 3D U-Net system for radiotherapy
dose prediction spanning two photon imaging modalities (CT and MRI) for
a public grand challenge, and a physics-based extension to proton
pencil-beam scanning. Using Task 1's first full training run's
evaluation results as a diagnostic signal rather than a final verdict,
we identified, empirically verified, and corrected a specific failure
mode -- dose leakage outside the patient body -- via an additional loss
term that then transferred to Task 2 without modification, by virtue of
a shared codebase. We further described a proton pencil-beam encoder
derived from established beam physics (Bragg-Kleeman range, Bortfeld
straggling, Highland scattering), and a GPU memory architecture fix
required to make its training tractable at proton's substantially larger
voxel grid. We report Tasks 1 and 2 as complete, evaluated, and packaged
for submission, and Task 3 as an in-progress extension with training
under way but no validated held-out results yet, and Task 4 as not yet
started. We view both the diagnostic process behind the Task 1--2 fix
and the first-principles physics derivation behind the Task 3 encoder as
useful methodological contributions independent of final leaderboard
standing.

% =====================================================================
\section*{Acknowledgments}

The author used an AI coding and writing assistant (Claude, Anthropic)
as a supporting tool during this project, under the author's direction
and review throughout. All research design, experimental decisions, and
result verification were the author's own. Specific tool-assisted uses
requiring disclosure per the submission's AI-usage policy: (1) code
implementation support for the beam encoders described in
Sections~\ref{sec:photon-encoder} and~\ref{sec:proton-encoder}, carried
out to the author's specifications, including the coordinate-grid
memory optimization; (2) drafting support during the out-of-body dose
leakage investigation of Section~\ref{sec:loss}, where the diagnostic
hypothesis and the decision to verify it against real ground-truth data
before changing any code were the author's own; (3) editing assistance
on this manuscript's text, drafted from the author's project notes and
experiment logs. All reported experimental results
(Tables~\ref{tab:t1-fix}--\ref{tab:t2} and
Section~\ref{sec:results}'s narrative figures) were produced by running
the author's own training and evaluation code against the challenge's
official dataset and metrics; no results in this paper were generated or
fabricated by the AI assistant.

% =====================================================================
\appendix
\section*{Artifact Description (AD) Appendix}

\textbf{Computational environment.} All training and evaluation was run
on a single NVIDIA A100 (40\,GB) GPU, with model code implemented in
PyTorch. Container packaging targets the Grand Challenge platform's
documented reference hardware (an AWS g5 instance).

\textbf{Software.} PyTorch (mixed-precision training via
\texttt{torch.amp}), SimpleITK (medical image I/O and resampling),
Docker (submission packaging, following the challenge's required
\texttt{invoke} HTTP API contract).

\textbf{Data.} The challenge's official training data
(photon: 77 patients $\times$ 3 beams $\times$ 180 control points;
proton: $\sim$77 patients, $\sim$81{,}000 beamlets), obtained from the
challenge organizers under the challenge's data-use terms. This data is
not independently redistributable by the authors.

\textbf{Availability.} Source code (training pipeline, beam encoders,
loss functions, evaluation harness, and Docker packaging scripts) is
maintained in a private repository at submission time and has not yet
been made publicly available; the authors intend to release it
following the challenge's submission deadline, subject to any
challenge-specific terms governing code release. No public artifact
identifier (DOI, public repository URL) is available to include in this
submission for that reason.

\textbf{Reproducibility notes.} Exact hyperparameters (learning rate,
loss term weights, batch/accumulation configuration, epoch counts per
task) are given in Sections~III-E–III-F and IV; the out-of-body loss fix
(Section~III-D) and the shared-coordinate-grid memory fix
(Section~III-C-1) are both described with enough detail to reimplement
independently of our specific codebase.

% =====================================================================
\begin{thebibliography}{9}

\bibitem{geant4}
S. Agostinelli \emph{et al.}, ``Geant4 -- a simulation toolkit,''
\emph{Nucl. Instrum. Methods Phys. Res. A}, vol. 506, no. 3,
pp. 250--303, 2003, doi: 10.1016/S0168-9002(03)01368-8.

\bibitem{unet3d}
{\"O}. {\c C}i{\c c}ek, A. Abdulkadir, S. S. Lienkamp, T. Brox, and
O. Ronneberger, ``3D U-Net: learning dense volumetric segmentation from
sparse annotation,'' in \emph{Proc. MICCAI}, 2016, pp. 424--432,
doi: 10.1007/978-3-319-46723-8\_49.

\bibitem{highland1975}
V. L. Highland, ``Some practical remarks on multiple scattering,''
\emph{Nucl. Instrum. Methods}, vol. 129, no. 2, pp. 497--499, 1975,
doi: 10.1016/0029-554X(75)90743-0.

\bibitem{bortfeld1997}
T. Bortfeld, ``An analytical approximation of the Bragg curve for
therapeutic proton beams,'' \emph{Med. Phys.}, vol. 24, no. 12,
pp. 2024--2033, 1997, doi: 10.1118/1.598116.

\bibitem{monai}
M. J. Cardoso \emph{et al.}, ``MONAI: An open-source framework for deep
learning in healthcare,'' \emph{arXiv:2211.02701}, 2022,
doi: 10.48550/arXiv.2211.02701.

\end{thebibliography}

\end{document}
